\section{Provenance des résultats et reproduction de l'extension}
\label{ap:procedencia-reproduccion}
\begingroup
\setlength{\parskip}{1pt}

L'exposé part du noyau APP--TRIT--TPK et maintient sa construction
corrélative du continu. Les lectures régionales, le produit natif et les
horloges irréductibles précèdent leurs réalisations analytiques. L'intégration
du centre, de la frontière et de la mémoire réunit des résultats du corpus HMT--MD avec
les compositions développées pour le double cercle. La provenance est
enregistrée par résultat, de sorte qu'une formulation réunie ou un certificat
supplémentaire ne rend pas rétrospectivement nouvelle la structure qu'il utilise.

\subsection{Dépendances réunies dans le corps du texte}

Le noyau comprend les définitions et les opérations communes, avec
ses deux figures. Les définitions, énoncés et preuves des cinq
développements précédents sont conservés intégralement.

La coordinatisation centrale, sa frontière relevée et les signatures de transport sont développées dans la section~\ref{sec:centro-frontera}. La distinction entre le point central \((5,5)\), la signature produit \(555555\) et les feuillets que son transport distingue est conservée dans la définition~\ref{def:firma555} et la proposition~\ref{prop:firma555}. Le calendrier de capacité et ses deux niveaux de retour sont démontrés dans les théorèmes~\ref{vac-add-teorema-primario} et~\ref{vac-add-teorema-secundario}. La géométrie de la coordinatisation n'est pas identifiée à l'historique complet : sa récupération utilise la loi multiaffine et les quotients conservés à la frontière.

Le transport du double cercle et les moments de mémoire sont réunis dans
la section~\ref{sec:transporte-momentos}. Les lecteurs de queue gamma de
la section~\ref{sec:lectores-gamma} reconstruisent les colonnes avec leurs
domaines et produits croisés. La section~\ref{sec:schur-coercividad}
établit l'estimation dans une fenêtre complète et son transport vers
des sous-espaces avec queues. Les égalités de conservation et la positivité
d'une forme sont des affirmations distinctes : chaque résultat de cette extension
conserve le domaine dans lequel il est démontré.

Les antécédents formels nécessaires à ces compositions sont
développés dans les annexes A--C : règles d'émission et histoires compatibles,
lectures régionales et retour nonadique, et registre dodécaphasique réversible.
Les équations et les preuves utilisées restent ainsi exposées avec
les domaines et les données de représentation requis par chaque résultat.

\clearpage
\subsection{Ordre de reproduction et portée des contrôles}

L'édition conserve les antécédents, les cinq développements précédents et les contrôles ciblés qui accompagnent leurs identités. Sa compilation atteste une réalisation matérielle lisible du texte ; la conservation documentaire et les vérifications informatiques ne certifient pas par extension tous les résultats d'un projet. L'apport de l'extension consiste à incorporer les compositions suivantes avec leurs démonstrations.

Le dossier qui accompagne le PDF conserve les sources \LaTeX{}, les cinq
développements en Markdown, les programmes et les reçus JSON. La reproduction
suit la dépendance mathématique des objets :

\begin{enumerate}
\item Construire les lecteurs de centre et de frontière en conservant le résidu et
le quotient, et vérifier la partition des 324 germes sous les transports
d'avancée et de demi-tour.
\item Évaluer les parties entières du calendrier des vacances au moyen
d'encadrements rationnels. Les bornes avec reste précèdent les chiffres et fixent
la certification de chaque position.
\item Vérifier les bilans de couture, les terminaux de mémoire et les
factorisations des moments ; conserver les termes croisés lors de la formation
des matrices conjointes.
\item Évaluer les bornes locales de queue gamma et du complément de Schur,
avec les domaines des estimations imprimés dans leurs énoncés.
\item Compiler le manuscrit et confronter les sources communes, les équations transférées, les références et la disposition des pages.
\end{enumerate}

Les contrôles rationnels et par intervalles vérifient les identités et
les bornes finies qu'ils déclarent. Les démonstrations fonctionnelles établissent
les affirmations universelles correspondantes. La conservation des
fichiers et une compilation correcte sont des contrôles éditoriaux séparés ;
aucun n'est présenté comme un substitut à une démonstration globale.

L'exécution ordinaire et optimisée des programmes utilise des conditions explicites : désactiver les assertions de l'interpréteur n'altère pas le critère d'acceptation. Le fichier d'instructions de la livraison spécifie les commandes et les dépendances de l'environnement. Les originaux précédents sont conservés séparément et ne sont pas écrasés par cette extension.
\par\endgroup
